\begin{center}
\bfseries \fontsize{24}{12} \selectfont Abstract
\end{center}
\vspace*{1\baselineskip}
Researchers and engineers managing modern distributed computing fabrics -- Edge, Multi-Cloud, and Hybrid environments -- routinely face complex job scheduling, siloed telemetry, and unpredictable hardware failures. Evaluating a new, intelligent scheduling policy directly on production infrastructure is risky and expensive, and is prone to catastrophic failure during link outages or thermal events. \textbf{RAP twin} addresses this by providing a high-fidelity digital-twin and planning sandbox that automates and optimizes job placement, scheduling, and performance analytics for distributed systems without ever touching physical infrastructure. The platform unifies static hardware descriptors (YAML) with streaming live telemetry through a bidirectional REST API, maintaining a single synchronized state machine (\texttt{dt.state.DTState}). Workloads are routed through a pluggable planning engine offering six strategies -- greedy, cheapest-energy, cheapest-cost, greenest, resilient/federated, and an RL-Markov policy -- each supporting a dry-run mode that predicts Job Completion Time, energy, and SLA risk without committing real resource locks. A deterministic Chaos Engine injects link outages, packet noise, and thermal derating into the simulated topology, while a \texttt{SelfHealingController} and \texttt{ResourceGuardian} automatically reroute at-risk jobs and reclaim abandoned reservations. Beyond the original five-module design, the platform now includes cost and carbon accounting for every plan, a chaos-as-CI resilience gate that fails a build on SLA or latency regression, a blast-radius finder that searches for the smallest fault set that breaks a job, exportable incident bundles, a cluster importer that derives a twin directly from a live Kubernetes cluster, and a what-if capacity planner. All of this is exposed through a Flask REST API (API-key gated on mutating routes) and a real-time React/TypeScript dashboard. Measured benchmarks on the project's own fabric show a concrete trade-off rather than a uniform win: the resilient planner buys fallback coverage at a latency cost (p95 1.72\,s vs.\ greedy's 1.26\,s, 83.3\% vs.\ 100\% SLA hit rate under a tightened deadline), while the RL-Markov planner shows an unresolved $\sim$75$\times$ cost anomaly relative to greedy that the project documents honestly as open future work rather than papering over. A suite of 129 automated tests and a live deployment validate that the platform is usable end to end. By offering an automated, scalable, and honestly-benchmarked emulation sandbox, RAP twin lets organizations statistically evaluate workload resilience and cost before committing a scheduling policy to expensive physical infrastructure.

\vspace*{0.5\baselineskip}
\noindent\textbf{Keywords:} Digital Twin, Distributed Systems, Task Scheduling, Chaos Engineering, Predictive Analytics, Monte-Carlo Simulation, Edge Computing, Reinforcement Learning, Cost and Carbon Accounting, Resilience Testing, Flask, React.
